\documentclass[11pt]{article}

\usepackage[T1]{fontenc}
\usepackage{lmodern}
\usepackage{amsmath,amssymb,amsthm}
\usepackage[margin=1in]{geometry}
\usepackage{enumitem}
\usepackage{xcolor}
\usepackage{microtype}
\usepackage[hidelinks]{hyperref}

\newcommand{\R}{\mathbb{R}}
\newcommand{\Om}{\Omega}
\newcommand{\norm}[1]{\left\lVert #1\right\rVert}
\newcommand{\Tmax}{T_{\max}}

\title{Paper 2, Theorem 1.2: Statement Amendment}
\author{Internal research note for the \texttt{ShenWork} formalization}
\date{12 July 2026}

\begin{document}
\maketitle

\begin{abstract}
The parameter range printed in Theorem 1.2 of Chen--Ruau--Shen permits the
mixed undamped branch $a>0$, $b=0$.  In that branch a positive spatially
constant solution grows exponentially, contradicting the asserted uniform
boundedness.  This note records the counterexample, locates the missing
hypothesis in the published proof, gives two corrected formulations, and
states the corresponding amendment needed in the Lean headline.
\end{abstract}

\section{Source checked}

The source examined here is the original submission, not a repository
paraphrase:

\begin{quote}
Le Chen, Ian Ruau, and Wenxian Shen, \emph{Chemotaxis models with
signal-dependent sensitivity and a logistic-type source, I: Boundedness and
global existence}, arXiv:2512.14858v1, 16 December 2025.
\end{quote}

The local archival copies are \path{paper2.pdf} and the submitted TeX source
\path{.tmp/arxiv/2512.14858.tar}.  As of 12 July 2026, arXiv lists no later
revision.  Theorem 1.2 assumes
$a,b\geq 0$ and $\beta\geq 1$.  It claims boundedness when $0<m<1$, and
boundedness plus global existence when $m=1$ and
\begin{equation}
  \chi_0 < \frac{2(2\beta-1)}{\max\{2,\gamma N\}}.
  \label{eq:chi-threshold}
\end{equation}
Remark 1.4(2) also says that $a$ and $b$ may vanish.  No preceding convention
excludes the mixed case $a>0$, $b=0$.

\section{Counterexample}

Let $a>0$, $b=0$, and let the initial datum be the positive constant
$u_0(x)=c$.  On any smooth bounded Neumann domain, set
\begin{equation}
  u(t,x)=c e^{at},
  \qquad
  v(t,x)=\frac{\nu}{\mu}\,u(t,x)^\gamma.
  \label{eq:constant-solution}
\end{equation}
Both functions are spatially constant.  Therefore all spatial derivatives and
the chemotaxis divergence vanish.  The elliptic equation for $v$ holds at
every time, while the parabolic equation reduces to
\begin{equation}
  u_t=a u.
\end{equation}
Thus \eqref{eq:constant-solution} is a positive global classical solution for
every $m>0$, every $\beta\geq 0$, and every $\chi_0\in\R$.  However,
\begin{equation}
  \norm{u(t,\cdot)}_\infty = c e^{at}\longrightarrow\infty.
\end{equation}
Consequently, both parts of the printed theorem fail in the branch
$a>0$, $b=0$.  Part (2) already fails for $m=1$, $\beta=1$, and $\chi_0=0$,
which satisfy \eqref{eq:chi-threshold}.

There is also a solution-independent formulation of the contradiction.  For
every positive classical Neumann solution with $b=0$, integration of the
parabolic equation gives
\begin{equation}
  M'(t)=aM(t),
  \qquad
  M(t):=\int_\Om u(t,x)\,dx.
  \label{eq:mass-ode}
\end{equation}
Positivity gives $M(t)>0$, so $a>0$ forces unbounded mass.  On a bounded
unit-volume domain,
\begin{equation}
  M(t)\leq \norm{u(t,\cdot)}_\infty.
\end{equation}
Hence an eventually bounded positive global solution cannot exist in this
branch, independently of uniqueness or of a preferred explicit solution.

\section{Where the proof loses the hypothesis}

The missing parameter restriction is visible in both parts of Section 4.

\begin{enumerate}[leftmargin=2em]
\item Proposition 2.4 supplies a uniform mass bound only in the two cases
\begin{equation}
  a=b=0,
  \qquad\text{or}\qquad
  a>0\ \text{and}\ b>0.
  \label{eq:prop24-cases}
\end{equation}

\item In Section 4.2, the proof of Theorem 1.2(1) drops the nonpositive term
$-b\int_\Om u^{p+\alpha}$ and later invokes Lemma 4.1 together with
Proposition 2.4 to replace a mass-dependent quantity by a time-independent
constant.  Proposition 2.4 supplies no such constant when $a>0$, $b=0$.

\item In Section 4.3, the proof of Theorem 1.2(2) reduces to $b=0$ and obtains
an inequality of the form
\begin{equation}
  \frac1p\frac{d}{dt}\int_\Om u^p
  \leq -\int_\Om u^p
       +C\left(\int_\Om u\right)^p.
  \label{eq:lp-forcing}
\end{equation}
It then claims a uniform-in-time $L^p$ bound by Gronwall.  That conclusion
requires a uniform mass bound.  In the omitted mixed branch,
\eqref{eq:mass-ode} makes the forcing term in \eqref{eq:lp-forcing} grow
exponentially, so Gronwall yields growth rather than a uniform bound.

\item The appeal to Theorem 1.1 for $\chi_0\leq0$ does not repair the mixed
case.  Theorem 1.1 itself treats only the two branches in
\eqref{eq:prop24-cases}.
\end{enumerate}

\section{Corrected parameter statement}

The exact correction directly supported by the written proof is
\begin{equation}
  (a=b=0)\quad\text{or}\quad(a>0\ \text{and}\ b>0).
  \label{eq:proof-supported}
\end{equation}
A slightly stronger and mathematically natural correction is
\begin{equation}
  a=0\quad\text{or}\quad b>0.
  \label{eq:maximal-correction}
\end{equation}
Under the standing assumptions $a,b\geq0$, condition
\eqref{eq:maximal-correction} is equivalent to excluding exactly
\begin{equation}
  a>0\quad\text{and}\quad b=0.
\end{equation}
The additional case $a=0$, $b>0$ needs only the elementary mass estimate
\begin{equation}
  M'(t)=-b\int_\Om u^{1+\alpha}\leq0,
\end{equation}
so it provides the same uniform mass input used in Section 4.

\medskip
\noindent
\fbox{%
\begin{minipage}{0.92\textwidth}
\textbf{Recommended amendment.}
Assume $\beta\geq1$ and either $a=0$ or $b>0$.  If $0<m<1$, then (1.23)
holds.  If $m=1$ and \eqref{eq:chi-threshold} holds, then (1.23) holds and
$\Tmax(u_0)=\infty$.
\end{minipage}}

\medskip
If the authors do not wish to add the elementary $a=0$, $b>0$ mass case, the
first sentence should instead assume \eqref{eq:proof-supported}.  Remark
1.4(2) should then say that both parameters may vanish simultaneously, rather
than suggesting all mixed nonnegative cases.

\section{Lean amendment}

The repository contains both a concrete and a parameter-general
interval-domain refutation in
\texttt{ShenWork/Paper2/IntervalDomainTheorem12Refutation.lean}:
\begin{verbatim}
not_Theorem_1_2_intervalDomain_when_a_pos_b_zero
not_Theorem_1_2_intervalDomain_of_a_pos_b_zero
\end{verbatim}
The proofs contain no \texttt{sorry} and depend only on
\texttt{propext}, \texttt{Classical.choice}, and \texttt{Quot.sound}.

The existing Lean definition has a second, independent fidelity problem in
part (1): it asks only for some bounded local solution on some positive finite
horizon.  The paper fixes the unique maximal solution and asserts boundedness
on its full maximal interval.  A faithful headline should use the existing
continuation objects in
\texttt{ShenWork/PDE/IntervalDomainExistence.lean}:
\begin{verbatim}
ReachableClassicalHorizon
ReachableArbitrarilyLong
FiniteContinuationAlternativeBranch
StandardContinuationAlternative
\end{verbatim}
A bare existential local witness does not formalize (1.23).  The amended
headline must distinguish a finite maximal-time alternative, available only
when reachable horizons are bounded above, from a genuine glued global branch
when they are unbounded.

Theorem 1.3 is unaffected by this particular $a>0$, $b=0$ counterexample,
because its original statement explicitly assumes $a,b>0$.  A separate
exponent-domain defect in the proof of Theorem 1.3(iv) is documented in
\path{paper2_theorem13_case_iv_amendment.pdf}.

\section*{Reference}

L. Chen, I. Ruau, and W. Shen, \emph{Chemotaxis models with signal-dependent
sensitivity and a logistic-type source, I: Boundedness and global existence},
arXiv:2512.14858v1 (2025),
\url{https://arxiv.org/abs/2512.14858}.

\end{document}
